\subsection*{Zadania: redukcja układu sił i momentów}

\subsubsection*{Zadanie 1}

Narysuj układ sił $F_i$ zaczepionych w punktach $A_i$. Wyznacz moment główny od tych sił w punkcie $C(3,3,0)\si{m}$.

\begin{equation*}
    \begin{cases}
        A_1(3;\ 5;\ 0) \\
        A_2(4.414;\ 4.414;\ 0) \\
        A_3(5;\ 3;\ 0) [ \si{m} ]
    \end{cases}
    \qquad
    \begin{cases}
        \Vec{F}_1 = [2;\ 0;\ 0 ] \\
        \Vec{F}_2 = [1.414;\ -1.414;\ 0 ] \\
        \Vec{F}_3 = [0;\ -2;\ 0 ] [ \si{kN} ]
    \end{cases}
\end{equation*}

\subsubsection*{Zadanie 2}

Narysuj układ sił $F_i$ zaczepionych w punktach $A_i$. Wyznacz moment główny od tych sił w punkcie $C(4,2,0)\si{m}$.

\begin{equation*}
    \begin{cases}
        A_1(0;\ 4;\ 0) \\
        A_2(4;\ 4;\ 0) \\
        A_3(8;\ 4;\ 0) [ \si{m} ]
    \end{cases}
    \qquad
    \Vec{F}_1 = \Vec{F}_2 = \Vec{F}_3 = [2;\ 0;\ 0] \si{kN}
\end{equation*}

\subsubsection*{Zadanie 3}

Narysuj układ sił $F_i$ i momentów $M_j$ zaczepionych w punktach $A_i,A_j$. Zredukuj układ do punktu $O(0,0,0)$.

\begin{equation*}
    \begin{cases}
        A_1(3;\ 2;\ 0) \\
        A_2(0;\ 3;\ 4) \\
        A_3(2;\ 0;\ 4) \\
        A_4( 0;\ 2;\ 4 ) \\
        A_5(0;\ 5;\ 0) [ \si{m} ]
    \end{cases}
    \qquad
    \begin{cases}
        \Vec{F}_1 = [ 5;\ 5;\ 0 ] \\
        \Vec{F}_2 = [ 10;\ 0;\ 0 ] \\
        \Vec{F}_3 = [ 0;\ -10;\ 0 ] \\
        \Vec{F}_4 = [-5;\ 0;\ 5] [ \si{kN} ] \\
        \Vec{M}_5 = [0;\ 0;\ 40] \si{kNm}
    \end{cases}
\end{equation*}